% SEGMENT OLP-0532-S01
% Part: set-theory
% Chapter: story
% Section: russells-paradox-again

\documentclass[../../../include/open-logic-section]{subfiles}

\begin{document}

\olfileid{sth}{story}{rus}
\olsection{ரஸ்ஸலின் முரண்பாடு (மீண்டும்)}

\olref[sfr][][]{part} இல் முறைசாரா கணக் கோட்பாட்டைப் பயன்படுத்தினோம்.
ஆனால் \emph{மிகவும்} முறைசாராத ஒரு கருத்தாக்கத்தின்படி, கணங்கள் என்பவை
பயனிலைகளின் விரிவுநிலைகள் மட்டுமே. இந்த முறைசாராக் கருத்து பின்வரும்
கொள்கையைக் கோரும்:

\begin{defish}
  \emph{முறைசாராப் பண்புவழிக் கணமாக்கல்.} எந்த வாய்பாடு $\phi$ க்கும்
  $\Setabs{x}{\phi(x)}$ உள்ளது.
\end{defish}

இந்தக் கொள்கை கவர்ச்சிகரமானதாக இருந்தாலும், அது மெய்ப்பிக்கத்தக்கவாறு
முரணுடையது. இதை \olref[sfr][set][rus]{sec} இல் கண்டோம்; ஆனால் இந்த
முடிவு அவ்வளவு முக்கியமானதும் அவ்வளவு நேரடியானதுமாக இருப்பதால், அதை
மீண்டும் கூறுவது பயனுள்ளது. சொல்லுக்குச் சொல்.

\begin{thm}[ரஸ்ஸலின் முரண்பாடு]
$R = \Setabs{x}{x \notin x}$ என்ற கணம் இல்லை.
\end{thm}

\begin{proof}
$R = \Setabs{x}{x \notin x}$ இருந்தால்,
$R \in R$ ஆகும்போது, அப்போதும் அப்போது மட்டுமே $R \notin R$; இது ஒரு
முரண்.
\end{proof}

ரஸ்ஸல் இந்த முடிவை 1901 ஜூன் மாதத்தில் கண்டறிந்தார். (ஆனால் நாம் இப்போது
வழங்கிய வடிவத்தில் அவர் முரண்பாட்டை அமைக்கவில்லை; ஏனெனில்
\emph{Grundgesetze} நூலில் விளக்கப்பட்ட ஃப்ரேகேயின் கணக் கோட்பாட்டை அவர்
ஆராய்ந்து கொண்டிருந்தார். \olref[blv]{sec} இல் இதற்குத் திரும்புவோம்.)
ஃப்ரேகேயின் அமைப்பிலுள்ள முரணை விளக்கி, 1902 ஜூன் 16 அன்று ரஸ்ஸல்
அவருக்கு எழுதினார். அந்தக் கடிதப் பரிமாற்றத்திற்கும் சிறிது பின்னணிக்கும்
\citet[pp.~124--8]{Heijenoort1967} ஐப் பார்க்கவும்.

இந்த இரண்டு வரி நிறுவல் \emph{தூய தருக்கத்தின்} ஒரு முடிவு என்பதை
வலியுறுத்துவது பயனுள்ளது. $\Setabs{x}{x \notin x}$ என்ற நமது
குறியீட்டில் உறுப்புசார் சமத்துவம் என்ற (தருக்கமல்லாத?)\ அடிகோளை நாம்
மறைமுகமாகப் பயன்படுத்தினோம் என்பது உண்மையே; ஏனெனில்
$\Setabs{x}{\phi(x)}$ என்பது, அது இருந்தால், $\phi$ களின்
\emph{ஒன்றே ஒன்றான} (உறுப்புசார் சமத்துவத்தால்) கணத்தைக் குறிக்கிறது.
ஆனால் உறுப்புசார் சமத்துவத்தின் சாயலையும்கூட பின்வருமாறு முடிவைக்
கூறுவதன் மூலம் தவிர்க்கலாம்: \emph{தன்னுறுப்பில்லாத கணங்களை, அவற்றை
மட்டுமே, உறுப்புகளாகக் கொண்ட கணம் இல்லை}. மேலும் இதற்கும் கணங்களுக்கும்
பெரிதாகத் தொடர்பில்லை. ரஸ்ஸலே கவனித்தபடி, \emph{தம்மைத்தாமே சவரம்
செய்யாத ஆண்களை, அவர்களை மட்டுமே, சவரம் செய்யும் ஆண் யாருமில்லை} என்று
துல்லியமாக இதே போன்ற பகுத்தறிதல் உங்களை முடிவுசெய்ய வைக்கும். அல்லது:
\emph{தம்மைத்தாமே முகராத குட்டைநாய்களை, அவற்றை மட்டுமே, முகரும்
குட்டைநாய் எதுவும் இல்லை}.
இப்படியே தொடரலாம். திட்டவடிவில், முடிவின் வடிவம் இதுவே:
\[
\lnot \exists x \forall z(Rzx \liff \lnot R zz).
\]
இது முதல்வரிசைத் தருக்கத்தின் ஒரு தேற்றத் திட்டம் மட்டுமே. எனவே நமது
கணக் கோட்பாட்டில் சிறு மாற்றங்களைச் செய்வதால் மட்டும் ரஸ்ஸலின்
முரண்பாட்டைத் தவிர்க்க முடியாது; நாம் கணக் கோட்பாட்டுக்குச்
\emph{செல்வதற்கு} முன்பே அது எழுகிறது. நாம் (செவ்வியல்) முதல்வரிசைத்
தருக்கத்தைப் பயன்படுத்தப் போகிறோம் எனில், $R = \Setabs{x}{x\notin x}$
என்ற கணம் இல்லை என்பதை \emph{ஏற்றுக்கொண்டே} ஆக வேண்டும்.

இதன் விளைவு இதுதான். முறைசாராப் பண்புவழிக் கணமாக்கலை ஏற்றுக்கொண்டபடியே
முரணை \emph{தவிர்க்க} விரும்பினால், \emph{கணக் கோட்பாட்டில்} மட்டும் சிறு
மாற்றங்களைச் செய்ய முடியாது. அதற்குப் பதிலாக உங்கள் \emph{தருக்கத்தை}
முழுமையாக மாற்றியமைக்க வேண்டும்.

செவ்வியல் அல்லாத தருக்கங்களைக் கொண்ட கணக் கோட்பாடுகள் முன்வைக்கப்பட்டுள்ளன
என்பது உண்மை. ஆனால் அவை---மிக மென்மையாகச் சொன்னாலும்கூட---வழமையற்றவை.
ரஸ்ஸலின் முரண்பாட்டை ஒரு நேரடியான இல்லாமை நிறுவலாகக் கருதி, அதனுடன்
எவ்வாறு வாழ்வது என்று கற்றுக்கொள்வதே வழமையான அணுகுமுறை. நாமும் அந்த
அணுகுமுறையையே பின்பற்றுவோம்.

\end{document}
